\section{Bounded collision compensation and its covariance}
\label{sec:collision-covariance}
The discontinuity of the section can be treated on the collision clock
without making its indicator a multiplier on a distribution space. The
argument uses bounded variation in the initial collision coordinates,
smooth approximation, and the collision-space spectral splitting. It
will give a covariance for the actual return record by an exact stopping
identity, rather than by an assumed spectral expansion of an induced
operator.

Write $\eta_R=\one_{Y_R^*}$ and recall
\begin{equation}\label{eq:compensation-normalizations}
 c_* = \frac{91}{10000\pi},\qquad
 \bar\tau_R=\frac{\sqrt3/2-\pi R^2}{2R},\qquad
 \bar G_R=\left(0,0,c_*^{-1},\frac{\bar\tau_R}{c_*}\right).
\end{equation}
Here $\bar G_R=\int G_R\dd\nu_R^*$ is the mean of the actual induced
record, by the Kac identities. On the common collision cylinder put
\begin{equation}\label{eq:bounded-compensation}
 h_R=f_R-\bar G_R\eta_R,
 \qquad s_R=c_*^{-1}\eta_R,
 \qquad S_{m,R}u=\sum_{j=0}^{m-1}u\circ T_R^j.
\end{equation}
Thus $\int h_R\dd\nu=0$ and $s_R\nu=\nu_R^*$, with the latter
measure regarded as a probability on $M$. The four components of $h_R$
are bounded, although those of $G_R$ need not be.

\subsection{Variation of a single physical collision}
For a function on $\T\times(-1,1)$, $\mathrm{BV}$ denotes bounded
variation with respect to the flat coordinates $(\alpha,p)$: both
first distributional derivatives are finite signed measures. All
variation norms below include the $L^1(\nu)$ norm. They are not norms
of the density of a pushed-forward record.

\begin{lemma}[Uniform single-collision variation]\label{lem:physical-bv}
There is a constant $B$ such that
\[
 \|f_R\|_\infty+\|f_R\|_{\mathrm{BV}}
 +\|\eta_R\|_\infty+\|\eta_R\|_{\mathrm{BV}}
 +\|h_R\|_\infty+\|h_R\|_{\mathrm{BV}}
 +\|s_R\|_\infty+\|s_R\|_{\mathrm{BV}}\le B
 \qquad(R\in I).
\]
The norms of vector-valued functions mean the sum of the component norms.
\end{lemma}
\begin{proof}
The uniform horizon bounds all possible next-disk centers in a fixed
finite subset of the triangular lattice. Cover the angular circle by
finitely many bounded rational charts
$t=\tan((\alpha-\alpha_0)/2)$, with uniformly bounded chart derivatives
and inverse derivatives. In such a chart $n_\alpha$ and $t_\alpha$ are
rational functions of $t$. The outgoing velocity
$\sqrt{1-p^2}n_\alpha+pt_\alpha$ is semialgebraic in $(t,p)$. For each
candidate center $d$, the admissibility inequalities and the entry root
\[
 (d-Rn_\alpha)\cdot v
 -\sqrt{R^2-|d-Rn_\alpha|^2+((d-Rn_\alpha)\cdot v)^2}
\]
are semialgebraic in $(R,t,p)$. Taking the least admissible positive
entry root and its lattice label involves only finitely many comparisons.
Choose any fixed tie convention on the singular set. Thus the bounded
single-collision record, extended arbitrarily by zero where no regular
root is defined, is a bounded semialgebraic family on each chart.
Only these one-collision functions, not an arbitrary long itinerary,
are used in this assertion.

We spell out the uniformity implication. A bounded semialgebraic family
$u(\lambda,x)$, with $x$ in a bounded interval, has a uniform finite
number of monotonicity intervals in $x$. To see this, apply differentiable
cell decomposition to its graph with the parameters $\lambda$ ordered
first, and refine by the sign of its fiber derivative where defined.
Uniform finiteness bounds the number of remaining exceptional points
and fiber intervals. On each interval the function is monotone or
constant. These are the cell decomposition and monotonicity results of
\cite[Sections 2.1--2.2 and 6.2]{Coste}. In particular the number does
not depend on the parameter or on the other fixed coordinate of a slice.
If $|u|\le M_0$ and there are at most $b_0$ intervals and exceptional
points, its one-dimensional distributional variation is at most
$C b_0 M_0$, including the jumps.

Apply this statement first with $(R,p)$ as parameters and then with
$(R,t)$ as parameters. Integrating the two slice-variation bounds
against the remaining coordinate proves the bounds for the two
first distributional derivatives. Equivalently, integration by parts
on almost every slice bounds the action on every compactly supported
$C^1$ test function by its supremum norm. A fixed finite chart partition
transfers these bounds to $\alpha$; any jumps at the chart endpoints
have uniformly bounded size and length. This proves the assertion for
$f_R$, including the square-root behavior at grazing.

The section is a fixed finite union of coordinate rectangles with
moving centers. Its perimeter and the number of its sides are uniformly
bounded, so the same assertion for $\eta_R$ follows directly, without
requiring the centers to be semialgebraic in $R$. The explicit means in
\eqref{eq:compensation-normalizations} are bounded on $I$. The remaining
bounds now follow from \eqref{eq:bounded-compensation}.
\end{proof}

\begin{lemma}[Mean-preserving smoothing]\label{lem:bv-smoothing}
There are positive linear operators $\mathsf S_\delta$, $0<\delta<1/2$,
which preserve the integral with respect to $\nu$ and satisfy
\begin{equation}\label{eq:bv-smoothing-bounds}
 \begin{split}
 \|\mathsf S_\delta u\|_\infty&\le\|u\|_\infty,\qquad
 \|u-\mathsf S_\delta u\|_1\le C\delta\|u\|_{\mathrm{BV}},\\
 \|\mathsf S_\delta u\|_{\mathrm{BV}}&\le C\|u\|_{\mathrm{BV}},\qquad
 \|\mathsf S_\delta u\|_{C^2(\alpha,\varphi)}
       \le C\delta^{-2}\|u\|_\infty .
 \end{split}
\end{equation}
The last norm uses $p=\sin\varphi$, as in the collision spaces.
\end{lemma}
\begin{proof}
Reflect $u$ evenly across $p=1$ and $p=-1$ to a function periodic with
period four in $p$, retaining the angular periodicity. Convolve with a
nonnegative, even, smooth product kernel of mass one and scale $\delta$,
then restrict to $-1<p<1$. Reflection symmetry is preserved, so the
integral over this half of the extended cylinder remains the original
integral. The reflected function has variation bounded by a fixed
multiple of the original variation; its traces match at the reflection
boundaries. Translation of a bounded-variation function by a vector $a$
changes its $L^1$ norm by at most $|a|$ times its variation. This follows
first for smooth functions by the fundamental theorem of calculus,
then for distributional derivatives by convolution and passage to the
limit. Integrating this inequality against the kernel gives the $L^1$
error. Convolution contracts variation on the extended cylinder.
The $L^1$ norms of the first two derivatives of the kernel are
$O(\delta^{-1})$ and $O(\delta^{-2})$. These facts prove all the flat
coordinate estimates. Composition with $p=\sin\varphi$ preserves the
stated $C^2$ upper bound, since the first two derivatives of $\sin$
are bounded. No inverse coordinate change at grazing is used.
\end{proof}

\subsection{An exponentially summable collision covariance}
\begin{proposition}[Correlations of bounded-variation observables]
\label{prop:bv-correlation}
For every $B_0<\infty$ there are $C,\gamma>0$, uniform in $R$, such
that centered real functions $a,b$ with
$\|a\|_\infty+\|a\|_{\mathrm{BV}}+
\|b\|_\infty+\|b\|_{\mathrm{BV}}\le B_0$ satisfy
\begin{equation}\label{eq:bv-correlation}
 \left|\int a(b\circ T_R^k)\dd\nu\right|\le Ce^{-\gamma k}
 \qquad(k\ge0).
\end{equation}
If also $\|a\|_1\le C_0\delta$, then for $r=a$ and $b=a$,
\begin{equation}\label{eq:small-bv-variance}
 \int|S_{m,R}r|^2\dd\nu
       \le C m\delta(1+|\log\delta|).
\end{equation}
The second constant may depend on $C_0$ and $B_0$, but not on $m,R$.
\end{proposition}
\begin{proof}
For smooth centered $a,b$, apply Lemma~\ref{lem:collision-spectral-input}
to the distribution $a\nu$ and evaluate the result on $b$ using smooth
multiplication and the mass functional. It gives
\[
 \left|\int a(b\circ T_R^k)\dd\nu\right|
       \le C\vartheta^k\|a\|_{C^2}\|b\|_{C^2}.
\]
Smooth both bounded-variation functions at scale $\epsilon$. Invariance
of $\nu$, boundedness, and \eqref{eq:bv-smoothing-bounds} bound the
change in the correlation by $C\epsilon$. The smooth correlation is
bounded by $C\vartheta^k\epsilon^{-4}$. Choose
$\epsilon=\tfrac14\vartheta^{k/5}$ and enlarge $C$ for $k=0$.
This proves \eqref{eq:bv-correlation}. In particular the constants depend
only on the displayed variation and supremum bounds, not on a modulus
of continuity or on a chosen smoothing scale.

For the last assertion, boundedness gives
$|\int r(r\circ T_R^k)\dd\nu|\le C\delta$ at every lag, while
\eqref{eq:bv-correlation} gives the exponential bound. Hence
\[
 \sum_{k\ge0}\left|\int r(r\circ T_R^k)\dd\nu\right|
       \le C\sum_{k\ge0}\min\{\delta,e^{-\gamma k}\}
       \le C\delta(1+|\log\delta|).
\]
Expanding the square of the stationary sum proves
\eqref{eq:small-bv-variance}.
\end{proof}

Put $h_{R,\delta}=\mathsf S_\delta h_R$. Define, on the collision clock,
\begin{equation}\label{eq:collision-gk}
 \begin{split}
 \Gamma_R={}&\int h_R\otimes h_R\dd\nu\\
 &+\sum_{k=1}^{\infty}\int
 \bigl[h_R\otimes(h_R\circ T_R^k)
       +(h_R\circ T_R^k)\otimes h_R\bigr]\dd\nu .
 \end{split}
\end{equation}
Use the same formula with $h_{R,\delta}$ for $\Gamma_{R,\delta}$.

\begin{theorem}[A continuous collision Green--Kubo matrix]
\label{thm:collision-covariance}
The series \eqref{eq:collision-gk} and its smoothed versions converge
absolutely, with exponential tails uniform in $R$ and $\delta$. The
matrix $\Gamma_R$ is continuous and positive semidefinite on $I$, and
\begin{equation}\label{eq:collision-covariance-limit}
 \Gamma_R=\lim_{m\to\infty}\frac1m
   \int S_{m,R}h_R\otimes S_{m,R}h_R\dd\nu,
 \qquad
 \|\Gamma_{R,\delta}-\Gamma_R\|\le C\delta(1+|\log\delta|).
\end{equation}
The first limit is uniform in $R$.
\end{theorem}
\begin{proof}
Lemma~\ref{lem:physical-bv}, smoothing, and
Proposition~\ref{prop:bv-correlation} give the uniform exponential tails.
For each fixed $k$, the integrand in \eqref{eq:collision-gk} converges
almost everywhere as $R\to R_0$. Indeed, exclude the finitely many
collision singularity preimages and section boundaries encountered up
to time $k$ at $R_0$. On the remaining full-measure set the physical
roots and finite itinerary persist and the moving rectangles have
stable membership. All observables are uniformly bounded. Dominated
convergence proves continuity at each fixed lag; the uniform tail then
proves continuity of the infinite series. This argument establishes
long-time uniformity through the correlation estimate, not through a
fixed-record continuity constant.

Expansion of the square gives the finite covariance with lag weights
$1-k/m$. Absolute summability, uniformly in $R$, proves its limit;
indeed the error is $O(m^{-1})$ by the exponential bound. Positivity
follows from positivity of each finite covariance. At each lag the
change made by smoothing is at most $C\delta$, using invariance and the
$L^1$ approximation. Both the original and the smoothed lag correlations
are at most $Ce^{-\gamma k}$. Summing the minimum of these two bounds
proves the final estimate.
\end{proof}

\begin{remark}\label{rem:collision-gk-scope}
The sum in \eqref{eq:collision-gk} is for the bounded, compensated
collision observable $h_R$ under $T_R$. It is not the Green--Kubo sum
of the unbounded observable $G_R-\bar G_R$ under $F_R^*$. Its relation
to the Gaussian law of the latter will be proved by stopping the actual
collision sums. No bounded-variation estimate for a many-return coarea
density, and no second derivative of that density, follows from the
single-collision variation lemma.
\end{remark}
